% !TEX root = Final_Report.tex

\documentclass[12pt]{article}
\usepackage[a4paper, total={6.5in, 9.5in}]{geometry}
%\usepackage[utf8]{inputenc}

\usepackage{amsmath}
\usepackage{amssymb}
\usepackage{graphicx}
\usepackage[usenames, dvipsnames]{color}

\usepackage{minted}
%\usepackage{hyperref}

\setlength{\parskip}{1em}
\renewcommand{\baselinestretch}{1.2}

\DeclareMathOperator{\Tr}{Tr}
\DeclareMathOperator{\sech}{sech}

\begin{document}

\begin{titlepage} % Suppresses displaying the page number on the title page and the subsequent page counts as page 1
	\newcommand{\HRule}{\rule{\linewidth}{0.5mm}} % Defines a new command for horizontal lines, change thickness here
	
	\center % Centre everything on the page
	
	%------------------------------------------------
	%	Headings
	%------------------------------------------------
	
	\textsc{\LARGE University of Cambridge}\\[1.5cm] % Main heading such as the name of your university/college
	
	\textsc{\Large Part III Physics}\\[0.5cm] % Major heading such as course name
	
	\textsc{\large Final Project Report}\\[0.5cm] % Minor heading such as course title
	
	%------------------------------------------------
	%	Title
	%------------------------------------------------
	
	\HRule\\[0.4cm]
	
	{\huge\bfseries Mean-Field Study of \\Kondo Phase Diagram}\\[0.4cm] % Title of your document
	
	\HRule\\[1.5cm]
	
	%------------------------------------------------
	%	Author(s)
	%------------------------------------------------
	
	\begin{minipage}{0.4\textwidth}
		\begin{flushleft}
			\normalsize
			\textit{Candidate}\\
			Elis \textsc{Roberts}\\
			\vspace*{5mm}
			\dotfill\hspace*{5mm}
		\end{flushleft}
	\end{minipage}
	~
	\begin{minipage}{0.4\textwidth}
		\begin{flushright}
			\normalsize
			\textit{Supervisor}\\
			Claudio \textsc{Castelnovo}\\
			\vspace*{5mm}
			\hspace*{5mm}\dotfill
		\end{flushright}
	\end{minipage}
	
	%------------------------------------------------
	%	Date
	%------------------------------------------------
	
	\vfill\vfill\vfill % Position the date 3/4 down the remaining page
	
	{\large\today} % Date, change the \today to a set date if you want to be precise
	
	%------------------------------------------------
	%	Logo
	%------------------------------------------------
	
	\vfill\vfill
	\includegraphics[width=0.18\textwidth]{crest.eps}\\[1cm] % Include a department/university logo - this will require the graphicx package
	 
	%----------------------------------------------------------------------------------------
	
	\vfill % Push the date up 1/4 of the remaining page
	
\end{titlepage}

\title{
  Mean-Field Study of Kondo Phase Diagram \\
  \large Part III Project Report
}
\author{\textit{Supervisors:} Claudio Castelnovo \& Garry Goldstein}
\date{\today}

% Start making a plan of what kind of things you need to talk about

\maketitle

\begin{abstract}
Phase diagram predictions for the Kondo model from mean-field theory are known to disagree with exact results described by a \emph{Bethe ansatz} solution. In this project we apply a new variation of mean-field theory to this problem, which differs in the novel way that it implements positive constraints into the Lagrangian through the introduction of an auxiliary system. This implementation introduces a new free parameter, which has previously been shown to enable a better a prediction of the characteristic energy scale, the Kondo temperature, and other derived properties of the system at zero temperature. Here we show that this free parameter must necessarily be made temperature dependent if it is to reproduce basic features of the finite temperature behaviour, but that after making this allowance it is possible to increase the order of the otherwise second-order phase transition. Nevertheless, it is shown that this added freedom is still not enough to avoid a phase transition entirely, as is typical of other mean-field theories in general.
\end{abstract}
% Talk about the goals of the project
% Summarise a few finite temperature results

\section{Introduction}

Metallic systems with localised magnetic impurities have, over the years, been the subject of much research in condensed matter physics, falling under the broader branch of strongly correlated systems which are characterised by interactions being significant in comparison to the kinetic energy dispersion (bandwidth).

The enormous theoretical challenge of studying strongly correlated systems, with their broad ranges of energy scales, means that one often turns to effective models to describe the low energy behaviour, introducing strict constraints that arise after \textit{projecting out} higher energy terms. One such effective model, studied in this project, is the Kondo model. This describes conduction electrons coupled to a (single) localised spin at the origin as a model for the dilute magnetic impurities that are sometimes present in a metal.

The many rich features of the general Kondo problem have been widely studied, with some simpler formulations being amenable to an exact solution via Bethe ansatz techniques. Often times, however, it is necessary to employ approximate methods to obtain results in more general cases, one such method being mean-field theory.

Use of mean-field theory is far from ideal, however, since current formulations applied to the Kondo impurity model are known to give results in disagreement with the Bethe ansatz solution for a characteristic energy of the problem known as the Kondo temperature $ T_{K} $ and (by extension) the magnetic susceptibility at zero temperature. The heat capacity is also greatly underestimated by existing mean-field methods. Recently, a new mean-field approach has been proposed by Garry Goldstein, Claudio Castelnovo (supervising the project) and Claudio Chamon which has given improved estimates of these quantities, which may be a sign that this new variation is indeed an improvement over existing formulations.

One significant aspect of the Kondo model that existing mean-field formulations have thus far failed to capture properly is the crossover from a Kondo to a paramagnetic phase in the temperature-field phase diagram, instead predicting a phase transition as in Figure~\ref{fig:phase_diagram}. The primary aim of this project is therefore to extend this new formulation to finite temperature, specifically to the temperature-field phase diagram and see whether a phase transition or a crossover is predicted. If the predicted behaviour is found to align with that of the exact Bethe ansatz solution, then the case for this new mean-field approach as an alternative to existing formulations would be greatly strengthened.

\begin{figure}[ht]
  \centering
  \includegraphics[width=0.6\textwidth]{Figures/phase_diagram.pdf}
  \caption{A representative phase diagram obtained from another (Read-Newns) mean-field approach (with calculations found in \cite{ManyBodyPhysics}), which shows a distinct phase transition between two phases. The Kondo temperature $ T_K $ is used to make quantities dimensionless.}
  \label{fig:phase_diagram}
\end{figure}

The report is structured as follows: Section~\ref{sec:theory} introduces the theoretical ideas underpinning the project; Section~\ref{sec:finite_temperature} contains the main finite temperature investigation; Section~\ref{sec:finite_field} gives a brief account of the behaviour at finite magnetic field and Section~\ref{sec:conclusion} concludes with an overview of the main results of the project.

% Citation needed

% Say what each section is going to contain

% !TEX root = Final_Report.tex

\section{Theoretical Background} % Remember to fix the sigma subscripts
\label{sec:theory}

\subsection{Kondo Model}

The Kondo model made its first appearance in 1964 when theorists were attempting to explain puzzling experimental observations made 30 years earlier that certain metals containing magnetic impurities showed minima in their resistivity as a function of temperature. Jun Kondo proposed the Kondo model to describe a new scattering mechanism introduced by magnetic impurities which accounted for the functional form of the resistivity.

In its simplest single impurity flavour, the Kondo model has the following Hamiltonian: \begin{equation} H_{\text{Kondo}}=\sum_{\boldsymbol{k},\sigma}\epsilon_{\boldsymbol{k}} c_{\boldsymbol{k},\sigma}^{\dagger}c^{}_{\boldsymbol{k},\sigma}+J\vec{S}\cdot\vec{s}(0) , \label{eq:KondoHamiltonian}\end{equation} in which there is a kinetic energy contribution from the conduction electrons $ c^{}_{\boldsymbol{k}, \sigma}$ as well as a term describing an antiferromagnetic coupling between a spin localised at the origin and the spin density of conduction electrons at that point.\footnote{The spin density of conduction electrons is given by $$ \vec{s}(0) = \frac{1}{2} \sum_{\boldsymbol{k}, \boldsymbol{k'}} \sum_{\sigma, \sigma'} c^{\dagger}_{\boldsymbol{k}, \sigma} \vec{\tau}_{\sigma, \sigma'} c^{}_{\boldsymbol{k'}, \sigma'} ~, $$ where $ \vec{\tau}_{\sigma, \sigma'} $ is a vector of Pauli matrices.}
% Check the factor of 1/2 here


% Should probably include a discussion of the Kondo model and things we know about its behaviour
% Talk about Kondo resonance and the Kondo temperature
% Talk about how the Kondo model is a low temperature limit to the Anderson model
% Indeed one of the reasons for the Kondo model's historical importance is how it lead to the development of the renormalisation group

% Falls within class of strongly correlated systems so usual band theory may not be enough like what might be possible with semiconductors etc
% Strongly correlated: can't think of things in terms of one-particle excitations
% Great intellectual challenge and so new methods to be able to go some way in predicting new materials

\subsection{Mean-Field Theory}

The theoretical calculations of this project are framed in terms of the path integral, which is an approach to statistical mechanics reminiscent of Feynman's path integral formulation of quantum mechanics. The partition function can be written as the functional integral over fermionic paths: $$ Z = \Tr{e^{- \beta H}} = \int \mathcal{D} [c^\dagger, c]~e^{-\int_{0}^{\beta} \,d\tau~L}\,, $$ from which many properties of the system may then be derived.\footnote{One also replaces fermion operators with \emph{Grassman numbers} within the integral, which have the property of anti-commutation (among others).} Here, the equivalent \emph{action} involves integration of the Lagrangian $ L $ over an imaginary time $ \tau = i t / \hbar $ with an upper limit of $ \beta = \frac{1}{k_B T} $.

% Put the k_B footnote here

% Don't be as informal with the word 'tricks'
One then begins to employ many `tricks' to make the problem more tractable. One such trick that proves to be useful is the introduction of new boson operators, which sometimes (as we shall see) necessitates that hard constraints be applied in the form of Lagrange multipliers. The way that constraints are implemented into the Lagrangian is shown in Appendix~\ref{sec:Constraints}.

The essence of mean-field theory is that we avoid performing the actual functional integration by approximating the integral by its saddle point, a step also known as the stationary phase approximation. In making this approximation, we are essentially imposing self-consistency conditions on whatever fields now appear in $ L $, making them take on their mean values. Thankfully, these mean-field self-consistency equations are exactly what result from directly minimising the effective action with respect to the auxiliary fields, which is what this project will involve.

% Maybe talk about how mean-field theory can be used to study BCS theory etc
% Maybe mention extensions like Dynamical Mean-Field theory

\subsubsection{Kotliar-Ruckenstein Bosons}

\subsection{The Soft-Constraint Approach}

% Make the k into a vector?

One existing mean-field approach is that of Read and Newns \cite{ReadNewns}, which represents the $ \vec{S}_0 $ term of the Hamiltonian in Eq~\eqref{eq:KondoHamiltonian} in terms of slave fermions $ f^{}_{\sigma} $, with occupation numbers obeying the hard constraint $ \sum_{\sigma} f^{\dagger}_{\sigma} f^{}_{\sigma} = 1 \,. $ The new mean-field approach that is being proposed reformulates this constraint as: \begin{equation} (1 - n_{\uparrow} - n_{\downarrow})^2 = n_{\uparrow} n_{\downarrow} + (1 - n_{\uparrow})(1 - n_{\downarrow}) = 0 . \label{eq:HardConstraint}\end{equation} Implementing this into the Lagrangian is formally equivalent, but problematic within mean-field theory because one later imposes: \[ \langle (1 - n_{\uparrow} - n_{\downarrow})^2 \rangle = 0\,, \] which, for such a positive semi-definite operator, enforces the exact constraint and leads to a diverging mean-field parameter. (Appendix~\ref{sec:Divergence} gives some feeling for why this is the case.)

A resolution to this issue is to begin by introducing an auxiliary non-interacting fermion $ h $ that is constrained to be trivially empty through imposing $ h^{\dagger} h = 0 $. One then combines this constraint with that of Eq~\eqref{eq:HardConstraint} by imposing instead: \begin{equation} n_{\uparrow} n_{\downarrow} + (1 - n_{\uparrow})(1 - n_{\downarrow}) - K h^{\dagger} h = 0\,, \end{equation} such that $ K > 0, K \neq 1 $ now encapsulates both constraints. Notice that this has introduced an arbitrary parameter $ K $ into the problem and thus a new degree of freedom, but has allowed us to circumvent the issues related to the previous hard constraint within mean-field theory. For this reason, this approach to mean-field theory has been internally referred to as the \emph{soft constraint approach}.

Proceeding in a similar fashion to the Read-Newns approach and incorporating four Kotliar-Ruckenstein (KR) \cite{KotliarRuckenstein} slave bosons, one obtains the following Lagrangian:
\begin{align}
\begin{split}
L \quad = \quad &\sum_{\boldsymbol{k},\sigma} c^{\dagger}_{\boldsymbol{k},\sigma} \left( \frac{d}{\,d\tau} + \epsilon^{}_{\boldsymbol{k}} - \mu \right) c^{}_{\boldsymbol{k},\sigma} + \sum_{\sigma} f^{\dagger}_{\sigma} \frac{d}{\,d\tau} f^{}_{\sigma} + h^{\dagger} \frac{d}{\,d\tau}h \\
&+ e^{\dagger} \frac{d}{\,d\tau} e  + \sum_{\sigma} p^{\dagger}_{\sigma} \frac{d}{\,d\tau} p^{}_{\sigma} + d^{\dagger} \frac{d}{\,d\tau} d\\
&+ \sum_{\sigma} \lambda^{}_{\sigma} (f^{\dagger}_{\sigma} f^{}_{\sigma} - p^{\dagger}_{\sigma} p^{}_{\sigma} - d^{\dagger} d^{} ) + \lambda_{\text{KR}} ( e^{\dagger} e + \sum_{\sigma} p^{\dagger}_{\sigma} p^{}_{\sigma} + d^{\dagger} d - 1)\\
&+\lambda_{\text{SC}} ( e^{\dagger} e + d^{\dagger} d - K h^{\dagger} h)\\
&+ 2 \frac{V V^{\ast}}{J} + \sum_{\boldsymbol{k},\sigma} \left( V^{\ast} c^{\dagger}_{\boldsymbol{k},\sigma} z^{}_{\sigma} f^{}_{\sigma} + V f^{\dagger}_{\sigma} z^{\dagger}_{\sigma} c^{}_{\boldsymbol{k},\sigma} \right) \,.\label{eq:Lagrangian}
\end{split}
\end{align}
The slave bosons $ e $, $ p_{\uparrow} $, $ p_{\downarrow} $ and $ d $ represent empty, singly occupied and doubly occupied states, respectively, as long as the fermion operator is suitably transformed to `update the books', so to speak:\footnote{The conventional transformation $ z^{}_{\sigma} \rightarrow (1 - d^{\dagger} d - p^{\dagger}_{\sigma} p^{}_{\sigma})^{- 1 / 2} z^{}_{\sigma} (1 - e^{\dagger} e - p^{\dagger}_{- \sigma} p^{}_{- \sigma})^{- 1 / 2} $ is also applied, but is not relevant to this discussion.} \begin{equation} f^{}_{\sigma} \rightarrow z^{}_{\sigma} f^{}_{\sigma}\,, \qquad z_{\sigma} = e^{\dagger} p^{}_{\sigma} + p^{\dagger}_{-\sigma} d \,. \end{equation} A Hubbard-Stratonovich transformation has also been applied to remove certain unpleasant terms in the Lagrangian, framing the interaction in terms of a new bosonic field $ V $.

% Don't say that the conventional transformation is not relevant, mention the slave boson review i.e. that including this means that the U=0 Anderson model can be recovered
% In this representation we have chosen to quantize along the z-axis, though it is possible to use Kotliar-Ruckenstein slave bosons in a rotationally invariant way

% From here, one may then derive the Helmholtz free energy and other quantities of interest within mean-field theory. The freedom in the parameter $ K $ allows for a whole family of mean-field solutions, from which one may obtain a single solution by tuning $ K $ to fit a known property of the Kondo model, for instance the Kondo temperature $ T_K $. Having tuned this parameter, other predictions about the system are made possible.


% !TEX root = Final_Report.tex

\section{Finite-Temperature Study}
\label{sec:finite_temperature}

As a first step towards constructing the temperature-field phase diagram, it is useful to gain some familiarity with the mean-field equations at finite temperature by investigating the Kondo model in the absence of a magnetic field. The isotropy of this zero-field case will allow for convenient simplification of some terms in the mean-field equations.

\subsection{Obtaining the Mean-Field Equations}
\label{subsec:obtaining_MF}

We start by deriving the self-consistency equations that must hold for the mean-field description of the system. Since $ F = - T \ln{Z} $, searching for the minimal action is equivalent to directly minimising of $ F $, illustrating a correspondence between the path integral and a more traditional way of approaching mean-field theory. Having introduced new bosonic fields to the Lagrangian of Eq~\eqref{eq:Lagrangian}, all fermionic fields may be integrated out as outlined in Appendix~\ref{sec:Free_Energy} to obtain an effective free energy
\begin{align}
\begin{split}
F \quad = \quad & \overbrace{- 2 T~\Re~{\sum_{\sigma}\ln{\left[\frac{\widetilde{\Gamma}(\xi_{\sigma} + D)}{\widetilde{\Gamma}(\xi_{\sigma})} \right]}}}^{F_0} + \frac{2 \Delta}{\pi \rho J} - \sum_{\sigma} \lambda_{\sigma} (p_{\sigma}^2 + d^2)\\
&+ \lambda_{\text{KR}} (e^2 + \sum_{\sigma} p_{\sigma}^2 + d^2 - 1) + \lambda_{\text{SC}} (e^2 + d^2) \underbrace{- T \ln{\left( 1 + e^{\beta K \lambda_{\text{SC}}} \right)}}_{F_{\text{h}}}
\label{eq:Free_Energy}
\end{split}
\end{align}

in terms of the \textit{gamma function} $ \widetilde{\Gamma} (z) \equiv \Gamma (\frac{1}{2} + \frac{z}{2 \pi i T}) $ and $ \xi_{\sigma} = \lambda_{\sigma} + i z_{\sigma}^2 \Delta $, a complex resonance-level energy made slightly different by the inclusion of the KR term. The intermediate summation that gives rise to $ F_0 $ has been regulated by a cut-off $ D $, the half-bandwidth, to reflect the fact that electrons of arbitrarily high energies do not exist within a metal.

% Give an explanation of terms

Note that the temperature dependence of this free energy is solely contained in $ F_0 $ and $ F_{\text{h}} $, which are the only terms that differ from the existing preliminary zero-temperature study of the soft-constraint approach. We now minimise this free energy to obtain a set of mean-field equations generalised to finite temperature, starting with the Hubbard-Stratonovich field

\begin{equation}
\frac{\partial F}{\partial \Delta} = 0 \implies \sum_{\sigma} \left[ \frac{\partial F_{0}}{\partial \xi_{\sigma}} - \frac{\partial F_{0}}{\partial \overline{\xi_{\sigma}}} \right] i z^2_{\sigma} = - \frac{2}{\pi J \rho} ~ .
\label{eq:MF_delta}
\end{equation}

Though strictly free to only treat bosonic variables as complex numbers, we also restrict our search to real solutions, leading to one equation for each KR boson\footnote{This may be partly justified by the phase invariance of most terms apart from $ z_{\sigma}^2 $, which is also why we have dropped modulus-squared signs for the KR bosons throughout.}: % Give more information in the footnote about how giving different phases to e d p could mess things up

\begin{align}
\frac{\partial F}{\partial d} = 0 &\implies \sum_{\sigma} \left[ \frac{\partial F_{0}}{\partial \xi_{\sigma}} - \frac{\partial F_{0}}{\partial \overline{\xi_{\sigma}}} \right] \frac{\partial z^2_{\sigma}}{\partial d} ~ i \Delta = -d ~ ( \lambda_{\text{KR}} + \lambda_{\text{SC}} -  \sum_{\sigma} \lambda_{\sigma} ) ~ , \label{eq:MF_d} \\
\frac{\partial F}{\partial e} = 0 &\implies \sum_{\sigma} \left[ \frac{\partial F_{0}}{\partial \xi_{\sigma}} - \frac{\partial F_{0}}{\partial \overline{\xi_{\sigma}}} \right] \frac{\partial z^2_{\sigma}}{\partial e} ~ i \Delta = -e ~ ( \lambda_{\text{KR}} + \lambda_{\text{SC}} ) ~ , \label{eq:MF_e} \\
\frac{\partial F}{\partial p_{\sigma}} = 0 &\implies \sum_{s} \left[ \frac{\partial F_{0}}{\partial \xi_{s}} - \frac{\partial F_{0}}{\partial \overline{\xi_{s}}} \right] \frac{\partial z^2_{s}}{\partial p_{\sigma}} ~ i \Delta = - p_{\sigma} ~ (\lambda_{\text{KR}} - \lambda_{\sigma}) ~ . \label{eq:MF_p}
\end{align}

The form of the $ \partial z^2_{\sigma} $ derivative terms are irrelevant for the time being, but are found in Appendix~\ref{sec:MF_eq_details}. Use of the chain rule means that the difficult derivative term
\begin{equation}
\frac{\partial F_{0}}{\partial \xi_{\sigma}} - \frac{\partial F_{0}}{\partial \overline{\xi_{\sigma}}} = \frac{i}{\pi} \Re{( \widetilde{\psi}(\xi_{\sigma} + D) - \widetilde{\psi}(\xi_{\sigma}))}
\label{eq:real_derivative}
\end{equation}
appears repeatedly, where $ \psi(z) \equiv \frac{d}{\,dz}(\ln{\Gamma(z)}) $ defines the \textit{digamma function}.

% Change all the - \sigma to a \overline{\sigma}

Finally, imposing Lagrange multiplier constraints completes the set of mean-field equations:
\begin{align}
\frac{\partial F}{\partial \lambda_{\sigma}} = 0 &\implies \left[ \frac{\partial F_{0}}{\partial \xi_{\sigma}} + \frac{\partial F_{0}}{\partial \overline{\xi_{\sigma}}} \right] = (p^2_{\sigma} + d^2) ~ , \label{eq:MF_lambda_sigma} \\
\frac{\partial F}{\partial \lambda_{\text{KR}}} = 0 &\implies e^2 + \sum_{\sigma} p^2_{\sigma} + d^2 = 1 ~ , \label{eq:MF_lambda_KR} \\
\frac{\partial F}{\partial \lambda_{\text{SC}}} = 0 &\implies e^2 + d^2 = K \langle h^{\dagger} h \rangle \equiv \kappa ~ , \label{eq:MF_lambda_SC}
\end{align}
where the other combination of difficult derivatives is
\begin{equation}
\frac{\partial F_{0}}{\partial \xi_{\sigma}} + \frac{\partial F_{0}}{\partial \overline{\xi_{\sigma}}} = - \frac{1}{\pi} \Im{( \widetilde{\psi}(\xi_{\sigma} + D) - \widetilde{\psi}(\xi_{\sigma}))} ~ .
\end{equation}

% Might want to write the derivative terms later

% Write down what the Helmholtz free energy is

% Go through some manipulation
% Explain how the mean occupation of the h comes in
% Argue that there is still freedom to choose the functional form of K as long as it's positive

\subsection{Solving the Mean-Field Equations}
\label{subsec:solving_MF}

Our aim is to self-consistently satisfy the mean-field equations derived above. The absence of any magnetic field means that there is nothing to favour any particular configuration of the spin-$\frac{1}{2}$ magnetic impurity, so we should expect the existence of a solution with $ p_{\uparrow} = p_{\downarrow} $ and $ \lambda_{\uparrow} = \lambda_{\downarrow} $, allowing us to drop spin indices $ \sigma $. Next, we shall make use of particle-hole symmetry to equate $ e^2 = d^2 $, which will dramatically simplify the mean-field equations.\footnote{Particle-hole symmetry comes out as a necessity in the Read-Newns mean-field approach, but here it is motivated by the belief that empty and doubly occupied pseudo-fermion states should be equally unphysical.}

% Want to have some discussion of particle hole symmetry

These simplifications mean that the occupation of each KR boson is entirely determined by $ \kappa $ as:
\begin{equation}
e^2 = d^2 = \frac{1}{2} ~ \kappa \quad \text{and} \quad p^2 = \frac{1}{2} (1 - \kappa) ~ .
\label{eq:soln_KR}
\end{equation}

Subtracting Eq~\eqref{eq:MF_e} and Eq~\eqref{eq:MF_d} immediately implies $ \lambda_{\sigma} = 0 $, which means that the complex energy $ \xi = i z^2 \Delta $ is now purely imaginary. Combining Eq~\eqref{eq:MF_lambda_sigma} and Eq~\eqref{eq:MF_lambda_KR} leads to $$ \frac{1}{2} = \frac{\partial F_0}{\partial \xi} + \frac{\partial F_0}{\partial \overline{\xi}} \approx - \frac{1}{\pi} \Im{\left[ \ln{\frac{D}{2 \pi i T}} - \widetilde{\psi}(\xi) \right]} ~ , $$ where we have used the large bandwidth $ D \gg T, \xi $ to make the leading order approximation that $ \psi(z) \approx \ln{z} $ for large $ z $. Taking the principal value of the complex logarithm then leads to the result that $$ \Im{\left[ \psi \left( \frac{1}{2} + \frac{\xi}{2 \pi i T} \right) \right] = 0} ~ , $$ which is merely consistent with the $ \lambda_{\sigma} = 0 $ conclusion that arose immediately from particle-hole symmetry and so tells us nothing new.

Turning to Eq~\eqref{eq:MF_delta} and making the same approximations, we may derive an implicit relation for $ \Delta $ in terms of $ T $ and $ \kappa $, similar to that found in \cite{ManyBodyPhysics}:
\begin{equation}
\psi \left( \frac{1}{2} + \frac{z^2 \Delta}{2 \pi T} \right) = \ln{\frac{D}{2 \pi T}} - \frac{1}{J \rho z^2} ~ .
\label{eq:soln_delta}
\end{equation}
Finding the finite temperature behaviour of the order parameter is therefore a case of inverting this relation for $ \Delta $, though it will not be possible to find an expression in terms of elementary functions. % Maybe citation

\subsubsection{Importance of $ \lambda_{\text{SC}} \geq 0 $}

In identifying $ \kappa = K \langle h^{\dagger} h \rangle $ as a free parameter, we should be certain that $ \lambda_{\textsc{SC}} > 0 $, since the thermal occupation of $ h $ takes on a familiar Fermi-Dirac form $$ \langle h^{\dagger} h \rangle = - \frac{1}{K} \frac{\partial F_{\text{h}}}{\partial \lambda_{\text{SC}}} = \frac{1}{1 + e^{- \beta K \lambda_{\text{SC}}}} ~. $$ If one were to find that $ \lambda_{\text{SC}} < 0 $, then the zero-temperature limit would have $ \langle h^{\dagger} h \rangle \rightarrow 0 $, thereby nullifying any effect of the freedom to choose $ K $.

Solving the remainder of the mean-field equations for $ \lambda_{\text{SC}} $, we find that
\begin{equation}
\lambda_{\text{SC}} = \frac{2 \Delta}{\pi J \rho} \frac{1 - 2 \kappa}{\kappa (1 - \kappa)} \geq 0 \quad \text{for} \quad \kappa \leq \frac{1}{2}
\label{eq:soln_lambda_SC}
\end{equation} which means $ \langle h^{\dagger} h \rangle \in \left(\frac{1}{2}, ~ 1\right) $ and so we may indeed treat $ \kappa $ as our free parameter (for the time being).\footnote{The added restriction on the magnitude of $ \kappa $ may be seen from Eq~\eqref{eq:soln_KR} to be equivalent to the seemingly reasonable statement that the physical pseudo-fermion states for the impurity should be more occupied than unphysical ones.}

\subsection{Zero-Temperature Heat Capacity}

As a quick check of validity of the above solution, one hopes that the zero-temperature limit should reproduce the leading order temperature dependence obtained from expanding the Fermi function around $ T = 0 $, something that has already been done using the soft-constraint approach \cite{Draft}.

Knowing that this particular response function is derived from the free energy as $ C = \frac{\partial E}{\partial T} = - T \frac{\partial^{2} F}{\partial T ^2} $, we start by finding the leading order behaviour of $ \frac{\partial F}{\partial T} $ from the minimised form of the free energy:
\begin{equation}
F^{\star} = \underbrace{- 4 T ~ \Re{ ~ \ln{\left[ \frac{\widetilde{\Gamma}(i z^2 \Delta + D)}{\widetilde{\Gamma}(i z^2 \Delta)} \right]}}}_{F_0^{\star}} + \frac{2 \Delta}{\pi J \rho} + \overbrace{\kappa \lambda_{\text{SC}}  - T \ln{\left( 1 + e^{\beta K \lambda_{\text{SC}}} \right)}}^{F_{\text{h}}^{\star}} ~ .
\end{equation}
It can be seen that the final term $ F_{\text{h}}^{\star} \approx - \kappa \lambda_{\text{SC}} e^{- \beta K \lambda_{\text{SC}}} $ will vanish quickly as $ T \rightarrow 0 $ in comparison to the first two terms, so we may safely neglect this contribution at leading order in $ T $.

Anticipating that explicitly calculating the inverse of Eq~\eqref{eq:soln_delta} will be somewhat difficult, we may find an expression for $ \frac{d \Delta}{d T} $ by inverting the chain rule
\begin{equation}
\frac{d \psi}{dT} = \frac{d \psi}{du} \frac{\partial u}{\partial \Delta} \frac{d \Delta}{d T} + \frac{d \psi}{du} \frac{\partial u}{\partial T} \implies \frac{d \Delta}{d T} = \left( \frac{d \psi}{dT} \left[ \frac{d \psi}{d u} \right]^{-1} - ~ \frac{\partial u}{\partial T} \right) \left[ \frac{\partial u}{\partial \Delta} \right]^{-1},
\end{equation}
where $ u \equiv \frac{1}{2} + \frac{z^2 \Delta}{2 \pi T} $ denotes the argument of $ \psi(u) $. Thus, the only non-trivial derivative term left to calculate is $ \left[ \frac{d \psi}{d u} \right]^{-1} $ which, using the asymptotic expansion of $ \ln{\Gamma(u)} $ \footnote{The asymptotic expansion of $ \ln{\Gamma(u)} $, known as \textit{Stirling's series}, is given by \cite{MathematicalFunctions} $$ \ln{\Gamma(u)} = \frac{1}{2} \ln{2 \pi} + u (\ln{u} - 1) - \frac{1}{2} \ln{u} + \frac{1}{12u} - \frac{1}{360 u^3} + \ldots ~. $$}, is $$ \left[ \frac{d \psi}{du} \right]^{-1} = ~ u - \frac{1}{2} + \frac{1}{12 u} + \ldots ~ . $$ This gives the leading-order temperature dependence of the order parameter as
\begin{equation}
\frac{d \Delta}{d T} \approx - \frac{\pi^2}{3} \frac{T}{z^4 \Delta} ~ ,
\label{eq:delta_derivative}
\end{equation}
reproducing a previous result in \cite{Draft}. Upon combining all derivative terms (as outlined in Appendix~\ref{sec:heat_capacity}), the zero-temperature heat capacity is found to be
\begin{equation}
C = \frac{2 \pi}{3} \frac{T}{z^2 \Delta}~ ,
\label{eq:heat_capacity}
\end{equation}
once again reassuringly consistent with the result obtained from the first-order correction to the Fermi function \cite{Draft}. The finite temperature mean-field equations therefore provide an alternative derivation of this limit of the heat capacity.

% Maybe put some of the calculations into an appendix?
% Mention that this had already been done in the draft paper by using a Fermi function approximation
% Say that it returns the same thing, which is a sign that no glaring errors have been made in solving the simultaneous equations
% Mention the Wilson ratio

\subsection{Plotting the Mean-Field Solution}

Satisfied that the zero-temperature limit of Eq~\eqref{eq:soln_delta} reproduces familiar results, we may use this implicit equation to plot the temperature dependence of the $ \Delta $ in the case of $ \kappa $ being constant\footnote{Here, $ \kappa $ is once again chosen such that $ z^{-2} = 1 - \frac{1}{2} \rho J \ln{(\rho J)} $ to reproduce $ T_K = D \sqrt{\rho J} e^{-1 / (\rho J)} $ \cite{Draft}.}, as shown in Figure~\ref{fig:delta_vs_T}.

\begin{figure}
\centering
\includegraphics[width=0.7\textwidth]{Figures/delta_vs_T.pdf}
\caption{A plot of the order parameter $ \Delta $ against the temperature $ T $, with both axes normalised by the Kondo temperature $ T_K $. The shaded region below the $ \Delta = 0 $ line indicates a non-physical region of the order parameter. The phase transition occurs at $ T_c \approx 1.13 ~ T_K $.}
\label{fig:delta_vs_T}
\end{figure}

From this plot it may be seen that the first derivative of the order parameter, $ \frac{d \Delta}{d T} $, is discontinuous at $ T = T_c $ if intervening by not allowing $ \Delta $ to become negative. This has the consequence (see Appendix~\ref{sec:discontinuity}) that $ \frac{d^2 F}{d T^2} $ is also discontinuous, characteristic of a second-order phase transition as opposed to a crossover. (Section~\ref{subsec:manipulating_PT} will investigate whether judiciously choosing the temperature dependence of the original soft-constraint parameter $ K \rightarrow K(T) $ can remove this phase transition.)

% Should maybe include the actual derivative here also

% Go through the work you did in trying to see if a phase transition can be removed
% Show the plot of the order parameter
% Want to know if the extra Fermi function dependence of kappa is still captured by K being a function of temperature

\subsection{Limitations of a Constant Soft-Constraint Parameter}

Before starting to artificially influence the temperature dependence of the model by promoting $ K \rightarrow K(T) $, one might question whether or not this is a natural thing to do in the first place. In solving the mean-field equations as we did in Section~\ref{subsec:solving_MF}, we absorbed the analytically difficult $ \langle h^{\dagger} h \rangle $ term into a free parameter $ \kappa $ that allowed us to reduce the mean-field equations to a single equation, namely Eq~\eqref{eq:soln_delta}. This strategy, however, came at the cost of expressing $ \Delta $ in terms of $ \kappa $ and \emph{not} the elementary parameter $ K $ which originally appeared in the introduction to the soft-constraint approach. We shall now observe the behaviour that would arise if insisting that $ K $ were instead held \emph{constant}.

Keeping the mean-field equations in terms of $ K $ and $ T $, we are left to solve two simultaneous equations involving $ \lambda_{SC} $ and $ \Delta $:
\begin{equation}
\psi \left( \frac{1}{2} + \frac{z^2 \Delta}{2 \pi T} \right) = \ln{\frac{T_K}{T}} - \ln{\left( 2 \pi \sqrt{\rho J} \right)} + \frac{1}{\rho J} \left[ 1 - \frac{1}{z^2} \right] ~ ,
\end{equation}

\begin{equation}
\quad \lambda_{\text{SC}} = \frac{8}{\pi \rho J} \frac{\Delta}{z^2} \left( 1 - \frac{2 K}{1 + e^{- \beta K \lambda_{\text{SC}}}} \right) ~ ,
\end{equation}

where $ z^2 = 4 K \left( 1 - K / (1 + e^{- \beta K \lambda_{\text{SC}}}) \right) / \left({1 + e^{- \beta K \lambda_{\text{SC}}}} \right) $. It may be appreciated that simply rearranging these equations for $ \Delta $ is now made impossible.

Nevertheless, it is still possible to make some progress by parameterising the above equations in terms of $ s \equiv \beta \lambda_{\text{SC}} $, and plotting the quantities $ \Delta / T_K $ and $ T / T_K $ along both axes according to:
\begin{equation}
\left( \frac{T}{T_K} \right) = \frac{1}{2 \pi \sqrt{\rho J}} ~ e^{(1 - z^{-2}) / (\rho J)} ~ \exp{\left[- \psi \left( \frac{1}{2} + \frac{\rho J ~ z^4 ~ s}{16 (1 - 2 K / (1 + e^{- K s}))} \right) \right]} ~ ,
\label{eq:parametric_T}
\end{equation}

\begin{equation}
\left( \frac{\Delta}{T_K} \right) = \frac{\pi \rho J ~ z^2}{8 (1 - 2 K / (1 + e^{- K s}))} \left( \frac{T}{T_K} \right) s ~ .
\label{eq:parametric_delta}
\end{equation}

\begin{figure}[ht]
\centering
\includegraphics[width=0.7\textwidth]{Figures/new_delta_vs_T_parametric.pdf}
\caption{A plot of the order parameter $ \Delta $ against the temperature $ T $, produced parametrically for the case of constant $ K $. The one-to-many mapping is not a physical function and so indicates a failure of the mean-field equations to pick out the minimal action.}
\label{fig:parametric}
\end{figure}

The resulting parametric plot, Figure~\ref{fig:parametric}, has a catastrophic shape in which there are no mean-field solutions beyond a certain temperature, but two extremal values of $ \Delta $ below this temperature.\footnote{This strange behaviour may look suspiciously like an artefact of the algebraic manipulations used to achieve the parametric form of Eq~\eqref{eq:parametric_T} and Eq~\eqref{eq:parametric_delta} (such as $ s $ not being an injective parameterisation), but this functional form is also reproduced when numerically solving the mean-field equations.} This suggests a breakdown of the mean-field equations and from this clearly unphysical behaviour we conclude that $ K $ must indeed have some temperature dependence if it is to accurately describe any finite temperature features of the model at all.

\subsection{Manipulating the Nature of the Phase Transition}
\label{subsec:manipulating_PT}

We have seen that neither constant $ K $ nor $ \kappa $ can reproduce a crossover from a strongly-coupled to weakly-coupled phase at finite temperature. We now investigate whether a suitable choice of $ \kappa \rightarrow \kappa(T) $ can remove this second-order phase-transition.

The first question to ask is whether it is possible for mean-field $ \Delta $ to only \emph{asymptotically} approach the weakly-coupled phase as temperature is increased, thereby not assuming a piecewise form. Consider choosing a different value for the soft-constraint parameter $ \widetilde{\kappa} $ such that $ z^2 \rightarrow \widetilde{z^2} = \alpha z^2 $. To see whether the form of $ \Delta(T) $ will change significantly upon making this choice, one can look at the new mean-field equation for $ \Delta $:
\begin{equation}
\psi \left( \frac{1}{2} + \frac{\alpha z^2 \Delta}{2 \pi T} \right) = \ln{\frac{D}{2 \pi T}} - \frac{1}{J \rho ~ \alpha z^2} ~ .
\end{equation}
When rescaling axes according to
\begin{equation}
\widetilde{T} = T \cdot e^{- \left( \frac{1}{\alpha} - 1 \right) / (J \rho z^2)} ~, \quad \widetilde{\Delta} = \Delta \cdot \alpha e^{- \left( \frac{1}{\alpha} - 1 \right) / (J \rho z^2)} ~ ,
\end{equation}
this equation reduces to nothing more than the original mean-field equation of Eq~\eqref{eq:soln_delta}, only this time in terms of $ \widetilde{\Delta} $ and $ \widetilde{T} $.
Therefore, choosing a different constant value of $ \kappa $ cannot change the shape of $ \Delta(T) $. Consequently, even if we promote $ \kappa \rightarrow \kappa(T) $ \footnote{After all, any temperature-dependent $ \kappa(T) $ is just switching between different `constant $ \kappa $' curves.}, the above mean-field condition will have an unavoidable $ \Delta = 0 $ solution at some finite temperature below
\begin{equation}
T^{\star}_c = T_K \cdot e^{- \psi \left( \frac{1}{2} \right)} / \left( 2 \pi \sqrt{\rho J} \right)~.
\end{equation}

Another option is to attempt to match successive derivatives of $ \Delta $ at the transition such that $ F $ will be continuous in all its derivatives. Trying to match all derivatives of a piecewise function that is analytic on both sides is conceptually dubious, however. This can be argued by saying that if all successive derivatives \textit{were} chosen to match, then both sides would share a Taylor series expansion about that point and hence be described by the same function, without need for a piecewise definition. Unless resorting to non-analytic functions, increasing the order of the phase transition is the closest one could get to a crossover.\footnote{Empirically, all derivatives of $ \Delta $ can be made to vanish by choosing $ \frac{1}{J \rho z^2} = \frac{1}{J \rho z_0^2} - \ln{\frac{T}{T_c}} $, but this forces $ \Delta = 0 $ at all temperatures (clearly undesirable) and contradicts the existing limits on the magnitude of $ \kappa $.}

With this in mind, we shall see what can nevertheless be achieved by choosing the form of $ \kappa (T) $. For example, let $ \kappa_0 $ be the value of the soft-constraint parameter that correctly picks out the Kondo temperature as before.\footnote{$ \kappa_0 = \frac{1}{2} \left( 1 - \sqrt{- \frac{1}{2} \rho J \ln{\rho J}} \right) $} Constructing $ \kappa(T) $ as
\begin{equation}
\kappa(T) = \kappa_0 + \delta \cdot \left[ \sech{\left( \frac{T}{t_1} \right)} + \tanh{\left( \frac{T - T_c}{t_2} \right)} \right] ~,
\end{equation}
can allow the order parameter to more gradually approach $ \Delta = 0$, as shown in Figure~\ref{fig:delta_smooth}, provided that $ t_1 $, $ t_2 $ and $ \delta $ are suitably chosen. The choice of the hyperbolic tangent above is somewhat arbitrary, but illustrates the principle that switching between $ (\kappa_0 - \delta) $ and $ (\kappa_0 + \delta) $ near the transition temperature can reduce the gradient to zero at the transition. (The hyperbolic secant term, on the other hand, serves to recover some of the zero-temperature behaviour.) There likely exist other, better choices for $ \kappa(T) $, depending on the desired functional form of $ \Delta(T) $; the only requirement is that $ \frac{d \Delta}{d T} \rightarrow 0 $ as $ \Delta \rightarrow 0 $, which can be achieved by tuning the parameters of the function such that 
\begin{equation}
\frac{d \kappa}{d T} \Bigr|_{T = T_c} = \frac{4 J \rho}{T_c} \frac{\kappa^2_0 (1 - \kappa_0)^2}{(1 - 2 \kappa_0)} \quad \text{and} \quad \kappa(T_c) = \kappa_0 ~ .
\end{equation}

\begin{figure}
\centering
\includegraphics[width=0.7\textwidth]{Figures/range_delta_vs_T.pdf}
\caption{A plot where the form of $ \kappa(T) $ is chosen to increase the order of the transition. Here, $ \kappa(T) $ is given by Eq~\eqref{eq:kappa_temp} (with $ \delta \approx 0.018 $, $ t_1 \approx \frac{1}{5} T_K $ and $ t_2 \approx \frac{4}{17} T_K $ for $ \rho J = 0.16 $) and switches from one constant $ \kappa $ curve to another. This allows the order parameter to take on any shape bounded by the blue shaded region, whose width is determined by the magnitude of $ \delta $.}
\label{fig:delta_smooth}
\end{figure}

One glaring deficiency of this analysis is that it is still unclear how $ K(T) $ should be chosen in the first place, since it would in general require solving
\begin{equation}
\kappa(T) = \frac{K(T)}{1 + e^{- \beta \lambda_{\text{SC}} K(T)}} ~ ,
\label{eq:kappa_temp}
\end{equation}
for which the simple relation $ K(T) \approx 2 \kappa(T) $ only holds near $ T \approx T_c $. \footnote{For what it's worth, a solution can be approached through the infinite expression
\begin{equation}
K = \kappa  \left( 1 + \exp{\left( - \beta \lambda_{\text{SC}} \kappa \left( 1 + \exp{\left( - \beta \lambda_{\text{SC}} \kappa \left( 1 + \ldots \right) \right)} \right) \right)} \right).
\end{equation}}

% Want to say that maybe the PT can be removed?

% !TEX root = Final_Report.tex

\section{Finite Magnetic Field}
\label{sec:finite_field}

Having explored the finite-temperature predictions for the soft-constraint in the absence of a magnetic field, we now introduce a finite magnetic field $ B $ which breaks the previous isotropy of the problem. One can include the effects of the magnetic field by introducing a Zeeman term for the magnetic impurity: $$ F_{B} = - g \mu_{B} B \left( f^{\dagger}_{\uparrow} f^{}_{\uparrow} - f^{\dagger}_{\downarrow} f^{}_{\downarrow} \right) \cong - g \mu_{B} B \sum_{\sigma} \sigma p^2_{\sigma} ~ , $$ where $ g $ is the \textit{Land\'{e} g-factor} and the effect on the conduction electrons has been ignored.

\subsection{Effect on the Mean-Field Equations}

The inclusion of a magnetic field leaves the previous mean-field equations of Section~\ref{subsec:obtaining_MF} largely unchanged apart from the equation for the singly-occupied KR boson, which now becomes
\begin{equation}
\frac{\partial F}{\partial p_{\sigma}} = 0 \implies \sum_{s} \left[ \frac{\partial F_{0}}{\partial \xi_{s}} - \frac{\partial F_{0}}{\partial \overline{\xi_{s}}} \right] \frac{\partial z^2_{s}}{\partial p_{\sigma}} ~ i \Delta = - p_{\sigma} ~ (\lambda_{\text{KR}} - \lambda_{\sigma} - \sigma g \mu_{B} B) ~ .
\end{equation}
Without isotropy, however, solving the mean-field equations becomes appreciably more difficult, and so we shall start with the slightly simpler task of searching for the location of the phase boundary $ \Delta = 0 $ in the $ B-T $ plane because it allows particle-hole symmetry to be justifiably used again.
% Maybe add an appendix in which you show that particle-hole symmetry works at \Delta = 0

\subsection{Plotting the Phase Boundary}

Setting $ \Delta = 0 $, the mean-field Lagrange multipliers are found to be:
\begin{equation}
\lambda_{\text{KR}} = 0 ~, \quad \lambda_{\text{SC}} = 0 ~,
\quad \lambda_{\sigma} = - \sigma g \mu_{B} B ~.
\end{equation}
Proceeding with this fixed $ B $-field, we obtain the condition that the occupancy of the impurity is given by
\begin{equation}
p^2_{\sigma} = \frac{1}{2} ( 1 - \kappa ) + \sigma \frac{1}{\pi} \Im{\left[ \psi\left( \frac{1}{2} + \frac{g \mu_B B}{2 \pi i T} \right) \right]} ~,
\end{equation}
which is problematic because $ \Im{\left[ \psi\left( \frac{1}{2} + \frac{g \mu_B B}{2 \pi i T} \right) \right]} \rightarrow - \frac{\pi}{2} $ in the large $ B $ limit, such that $ p^2_{\downarrow} $ becomes negative. Physically, this corresponds to one spin-state of the impurity being heavily favoured compared to the other, but this mean-field solution requires intervention with a piecewise definition to ensure that both $ p^2_{\sigma} $ saturate sensibly.
% Stop using the word problematic so much

Making a parametric plot of the phase boundary in terms of $ b = \frac{g \mu_B B}{T} $ according to
\begin{equation}
\left( \frac{T}{T_K} \right) = \frac{1}{2 \pi \sqrt{\rho J}} e^{\left( 1 - z^{-2} \right) / (\rho J)} \exp{\left[ - \Re ~ \psi \left( \frac{1}{2} + \frac{b}{2 \pi i} \right) \right]} \quad \text{and} \quad \frac{g \mu_B B}{T_K} = \left( \frac{T}{T_K} \right) b
\label{eq:phase_parametric}
\end{equation}
reveals a more worrisome problem with this approach, namely that the behaviour is dominated by the $ \frac{1}{\rho J z^2} $ that appears in the exponential. If $ \kappa $ is left constant, then this leads to a phase boundary that curves upwards at large magnetic fields, as shown in Figure~\ref{fig:new_phase_boundary}, because $ z^2 $ includes other field-dependent terms.\footnote{This strange phase boundary does not arise if one is actually interested in the case where $ B $ is not fixed, but rather a uniformly distributed stochastic field $ \widetilde{B} \in [ - B / 2, ~ B / 2 ] $, since the imaginary part of $ \widetilde{\psi} (g \mu_B \widetilde{B}) $ averages to zero and the same large-N limit of the phase boundary is recovered for constant $ \kappa $.} This is despite the fact that the rest of the expression behaves similarly to the large-N limit found by Piers Coleman \cite{ManyBodyPhysics}, which was used to create the original phase diagram of Figure~\ref{fig:phase_diagram}.

\begin{figure}
\centering
\includegraphics[width=0.7\textwidth]{Figures/new_phase_diagram.pdf}
\caption{The new phase diagram for the soft-constraint approach when $ \kappa $ held constant, generated parametrically using Eq~\eqref{eq:phase_parametric}. The shape of the boundary is highly sensitive to $ z^2(\frac{B}{T}) $, which is itself made to saturate by-hand beyond a certain value of $ \frac{B}{T}$. It is this manual intervention that leads to the sharp kink in the boundary at low temperatures. The dotted line shows the original large-$N$ boundary of Figure~\ref{fig:phase_diagram}.}
\label{fig:new_phase_boundary}
\end{figure}
% Need to have the form of z^2 in the appendix

If one is truly interested in the case of the Kondo model subject to a fixed magnetic field, then it appears that $ \kappa $ would have to acquire a (likely complicated) $B$-field dependence to tame this unexpected behaviour. Given that it was only ever possible to increase the order of the phase transition in the absence of a magnetic field, it seems likely that a phase transition is also inevitable at finite fields. (Even the task of increasing the order of a phase transition is made much more challenging at finite field, and so has not been attempted in this project.)

% Go through your result for the phase-boundary when K is constant, then frame it as if constant K is once again usatisfactory

% Mention that you could have done something similar where you average over the stochastic field, which makes z^2 into a constant thing so that you get back to the large-N limit

% Say that you haven't worked out how you would go about changing k at finite field to get a sensible curve

% Maybe add an appendix in which you show that particle-hole symmetry works at \Delta = 0

% Mention the way that a magnetic field can be included in an ad-hoc way
% Show that it is problematic unless you define a stochastic field
% Or just make the p_\sigma s saturate at zero? - Doesn't work

\section{Conclusion}
\label{sec:conclusion}

This project aimed to explore the finite temperature and field behaviour of the Kondo model as predicted by a promising new variation of mean-field theory: the soft-constraint approach. Failure of conventional mean-field theories to accurately describe the phase diagram of this model makes it an ideal platform to test whether this new method is indeed a significant improvement that could find further use in describing strongly correlated systems.

The vast majority of effort was spent on the extension of this approach to finite-temperature, which was enough to draw some conclusions about the behaviour of the soft-constraint away from the $ (T = 0,~B = 0) $ point in the phase diagram. Crucially, it was found that the associated soft-constraint parameter  has to be allowed to vary with temperature if it is to lead to a solvable set of mean-field equations. This is because the thermal occupancy of the auxiliary fermions introduced is enough to have a significant effect on the \emph{effective} soft-constraint parameter.

Even when allowing the soft-constraint parameter to vary with temperature, a phase transition is found to be inevitable, in disagreement with the crossover obtained from other, more intensive methods. Nevertheless, it was demonstrated that artificially choosing the form of a temperature-dependent soft-constraint parameter can increase the order of the phase transition, though at the cost of zero-temperature accuracy.

The increased difficulty of the mean-field equations in the presence of a magnetic field meant that this aspect of the investigation was limited to a phase boundary calculation. These results suggest that the soft-constraint parameter would have to acquire further field dependencies to better describe the location and nature of the phase transition in the temperature-field plane.

\newpage

% !TEX root = Final_Report.tex

\appendix

\section{Constraints in the Lagrangian}
\label{sec:Constraints}

The way that constraints can be implemented into the Lagrangian is illustrated in the following example. Suppose that we wanted to implement the constraint $ \sum_{\sigma} f^{\dagger}_{\sigma} f^{}_{\sigma} = 1 $, say, which would be equivalent to having a partition function $$ Z = \Tr{\left[ e^{- \beta H}~\delta{\left( \sum_{\sigma} f^{\dagger}_{\sigma} f^{}_{\sigma} - 1 \right)} \right]} . $$ We could then express the constraint as $$ \delta{\left( \sum_{\sigma} f^{\dagger}_{\sigma} f^{}_{\sigma} - 1 \right)} = \int_{0}^{2 \pi} \frac{\,d\alpha}{2 \pi}~e^{- i \alpha (\sum_{\sigma} f^{\dagger}_{\sigma} f^{}_{\sigma} - 1)} = \int_{0}^{2 \pi i k_B T} \frac{\,d \lambda}{2 \pi i k_B T} e^{- \beta \lambda (\sum_{\sigma} f^{\dagger}_{\sigma} f^{}_{\sigma} - 1)}, $$ where we have written $ \lambda = i \alpha k_B T $. Absorbing various factors into the measure of integration, we may now write: $$ Z = \int \, \mathcal{D} [\lambda]~\Tr{\left[ e^{- \beta H} e^{- \beta \lambda (\sum_{\sigma} f^{\dagger}_{\sigma} f^{}_{\sigma} - 1)} \right]} \,. $$ Imposing this constraint can therefore be seen to be equivalent to modifying the original path integral and including an extra term in the Lagrangian: $$ L \rightarrow L - \lambda \left( \sum_{\sigma} f^{\dagger}_{\sigma} f^{}_{\sigma} - 1 \right) \,. $$ In fact, this is actually the Read-Newns constraint that is imposed on the occupation of the fermions $ f^{}_{\sigma} $ (with $ \sigma \in \{ \uparrow , \downarrow \}$) representing the localised spin of the magnetic impurity.

\section{Divergent Mean-Field Parameter}
\label{sec:Divergence}

To see why imposing $ \langle (1 - n_{\uparrow} - n_{\downarrow})^2 \rangle = 0 $ leads to a divergent mean-field parameter, one may appreciate that by virtue of positive semi-definiteness, the mean-field condition \[ \frac{\, \delta Z}{\, \delta \lambda(\tau)} \Bigr|_{\bar{\lambda}} = 0 \] essentially becomes a condition on the integrand itself (namely something like $ P e^{- \int \, d\tau \bar{\lambda} P} = 0 $ for the constraint $ P $), which forces $ \bar{\lambda} \rightarrow \infty $.

\section{Deriving the Helmholtz Free Energy}
\label{sec:Free_Energy}

The Lagrangian of Eq~\eqref{eq:Lagrangian} now has all fermionic fields in quadratic form, since this was the purpose of the auxiliary bosons in the first place. One may therefore perform standard Gaussian integration over the Grassman variables to get an expression involving the determinant of the action,

% TODO: Make this equation less of an eyesore

\[
L \quad = \quad \sum_{\sigma}
\begin{pmatrix} \cdots & c^{\dagger}_{k,\sigma} & \cdots & f^{\dagger}_{\sigma} \end{pmatrix}
\begin{pmatrix} & & & \\ &(\epsilon_{k} + \partial_{\tau}) \delta_{k, k'} & & V^{\ast} z^{}_{\sigma} \\ & & & \\ & V z^{\dagger}_{\sigma} & & (\lambda_{\sigma} + \partial_{\tau}) \end{pmatrix}
\begin{pmatrix} \vdots \\ c^{\dagger}_{k',\sigma} \\ \vdots \\ f^{\dagger}_{\sigma} \end{pmatrix}
\quad + \quad \hdots
\]

\[
\rightarrow \qquad \sum_{\sigma, n} \quad \ln \det
\begin{pmatrix} & & & \\ &(\epsilon_{k} - i \omega_{n}) \delta_{k, k'} & & V^{\ast} z^{}_{\sigma} \\ & & & \\ & V z^{\dagger}_{\sigma} & & (\lambda_{\sigma} - i \omega_{n}) \end{pmatrix}
\quad + \quad \hdots \quad ,
\]

which involves a summation over Matsubara frequencies $ \omega_{n} $. The mean-field impurity electron contribution to the Helmholtz free energy is therefore

\[
F = - T \sum_{\sigma, n} \ln{\left(- i \omega_{n} + \lambda_{\sigma} + \sum_{k} \frac{z_{\sigma}^{2} |V|^{2}}{i \omega_{n} - \epsilon_{k}}\right)} \quad + \quad \cdots \quad ,
\]

where $ \cdots $ now includes the conduction electron contribution and all the other constraint terms previously present in the Lagrangian.

% Pages 675, 679 & 730 (+ 593 for the sgn(omega) part)
% All it does is put in V*z instead of V

% Will have to involve some Matsubara frequency stuff
% Go through the derivation of the Helmholtz free energy for the question, which should be pretty similar to Piers Coleman's result
% Probably leave this until last because it isn't essential

\section{Manipulating the Mean-Field Equations}


\begin{thebibliography}{1}
% Maybe look at doing the bibliography with a separate .bib file

\bibitem{ManyBodyPhysics}
P. Coleman
\emph{Introduction to Many Body Physics}
Cambridge University Press (2015)

\bibitem{Draft}
G. Goldstein, C. Castelnovo and C. Chamon
\emph{Mean-Field Method for Handling Null Constraints on Positive-Semidefinite Operators}
Draft (2017)

\bibitem{KotliarRuckenstein}
G. Kotliar and A. E. Ruckenstein
\emph{New Functional Integral Approach to Strongly Correlated Fermi Systems: The Gutzwiller Approximation as a Saddle Point}
Phys. Rev. Lett. \textbf{57}, 1362 (1986)

\bibitem{ReadNewns}
N. Read, and D. M. Newns
\emph{A New Functional Integral Formalism for the Degenerate Anderson Model} Journal of Physics C: Solid State Physics 16.29 (1983)

\bibitem{MathematicalFunctions}
M. Abramowitz
\emph{Handbook of Mathematical Functions, With Formulas, Graphs, and Mathematical Tables}
Dover Publications, Inc. (1974)

\bibitem{BetheAnsatz}
N. Andrei, K. Furuya, and J. H. Lowenstein
\emph{Solution of the Kondo Problem}
Reviews of Modern Physics 55.2 (1983): 331-402

\bibitem{SlaveBosons}
Frésard R., Kroha J. and Wölfle P. (2012)
\emph{The Pseudoparticle Approach to Strongly Correlated Electron Systems}
Strongly Correlated Systems. Springer Series in Solid-State Sciences, vol 171.

% https://dl.acm.org/citation.cfm?id=1830479

% Include Piers Coleman's slave boson mean field theory paper
% Include large N
% Include the Slave boson review

% MAYBE USEFUL: New approach to the mixed-valence problem
%http://iopscience.iop.org/article/10.1088/0022-3719/19/17/017/pdf

\end{thebibliography}

\end{document}